\chapter{Lean and Mathlib}
\label{chap:introduction}

The practice of mathematics is fundamentally rooted in the rigorous verification of proofs.
However, as mathematical theories grow increasingly complex, the traditional peer-review process
becomes strained. The digitization and formal verification of mathematics offer a modern solution
to this problem, allowing mathematicians to verify their results with absolute certainty using
software.

This chapter introduces the Lean theorem prover and its mathematical library,
called \lstinline|Mathlib|.

\section{The Lean Theorem Prover}
\label{sec:lean_theorem_prover}

Lean is a functional programming language and interactive theorem prover initially developed
by Leonardo de Moura at Microsoft Research in 2013 \cite{demoura2015lean}.
The current version, Lean 4,
is a complete reimplementation that serves simultaneously as a proof assistant and a
general-purpose, practically efficient programming language.

At its logical core, Lean is based on a variant of dependent type theory known as the
Calculus of Inductive Constructions. In this framework, propositions are represented
as types, and proofs are represented as terms of those types—a correspondence known as the
Curry--Howard isomorphism.

To prove a theorem in Lean is to construct a term of the corresponding type. Lean's kernel,
a small and highly trusted piece of code, checks that this term is well-typed. If the kernel
accepts the term, the proof is logically sound. This minimizes the trusted code base; users
do not need to trust the complex tactics or automation used to generate the proof, only the
kernel that verifies the final proof term.

% Unlike traditional computer algebra systems (such as Mathematica or SageMath) which
% primarily perform numerical and symbolic computations, interactive theorem provers
% like Lean are designed to reason about abstract mathematical concepts, properties, and
% logical deductions.

\section{Dependent Type Theory}
\label{sec:dependent_type_theory}

Traditional set theory (such as Zermelo--Fraenkel with Choice, ZFC) has historically served
as the foundation of mathematics. However, dependent type theory provides a more natural
environment for machine verification.

In Lean's type theory, every expression has a specific type. The defining feature of
\emph{dependent type theory} is that types can depend on terms. For example,
the type of a vector of length $n$, denoted \lstinline|Vector α n|, depends on the natural
number $n$, where \lstinline|α| is a type. This allows for highly expressive typing where the properties of mathematical
objects are baked directly into their definitions.

% Lean distinguishes between two fundamental universes:
% \begin{itemize}
%     \item \texttt{Prop}: The universe of logical propositions. Types residing in \texttt{Prop} are
%     subject to proof irrelevance, meaning any two proofs of the same proposition are considered
%     definitionally equal.
%     \item \texttt{Type}: The universe of data and mathematical structures
%     (e.g., natural numbers, rings, manifolds).
% \end{itemize}

Moreover, even types themselves have types. The type of propositions \lstinline|Prop| is also
of type \lstinline|Type|. There is a \emph{hierarchy of
type} to prevent to not run into a paradox called \emph{Girard's paradox}.
Girard's paradox is an inconsistency that arises in certain very expressive type systems,
especially those allowing a ``type of all types.'' It is the
type-theoretic analogue of Russell's paradox in naive set theory. Roughly
speaking, if a system allows types to quantify over themselves too freely,
then one can construct a self-referential type that leads to a contradiction.
More discussion of this can be found in \cite{coquand1986analysis}.
% \cite{Coquand1986}.

In hierarchy of type, the type of \lstinline|Type| is \lstinline|Type 1|, the
type of \lstinline|Type 1| is \lstinline|Type 2|, and so forth. When we
want to declare that \lstinline|R| is a type residing in an arbitrary
universe within this hierarchy, we write \lstinline|R : Type*|.


\section{A brief introduction to Lean syntax}

Lean is both an interactive theorem prover and a functional programming language.
Definitions, statements, and proofs are written in a syntax that is
subsequently verified by the Lean kernel. In this section, we introduce several
foundational pieces of Lean syntax that appear frequently in later chapters.

\subsection{Basic declarations}

A typical Lean file consists of a series of declarations. Beyond the foundational
\lstinline|def|, \lstinline|theorem|, and \lstinline|axiom|, Lean provides
powerful organizational tools such as \lstinline|abbrev|, \lstinline|structure|,
\lstinline|class|, and \lstinline|instance|.

The keyword \lstinline|def| is used to define new data, functions, or values.
For example, the following code defines the successor of a natural number:

\begin{lstlisting}[language=Lean]
def succ (n : Nat) : Nat := n + 1
\end{lstlisting}

Here, \lstinline|succ| is the name; \lstinline|(n : Nat)| declares an explicit input \lstinline|n| of type \lstinline|Nat|; the final \lstinline|: Nat| gives the return type; and \lstinline|:=| introduces the defining expression on the right.

Closely related is the keyword \lstinline|abbrev|. It also creates a definition, but instructs Lean's elaborator to treat the new name as fully transparent during unification. It is frequently used for type aliases:

\begin{lstlisting}[language=Lean]
abbrev Vector3 := Float × Float × Float
\end{lstlisting}

Lean then automatically unfolds \lstinline|Vector3| into a triplet of floats during type checking, avoiding the strict type-mismatch errors that a standard \lstinline|def| might cause.

The keyword \lstinline|theorem| is reserved for stating and proving propositions. A theorem requires a name, a logical statement, and a formal proof. For example:

\begin{lstlisting}[language=Lean]
theorem add_zero (a : Nat) : a + 0 = a := by
  rw [Nat.add_zero]
\end{lstlisting}

This theorem asserts that for every natural number \lstinline|a|, the equality \lstinline|a + 0 = a| holds. The portion before \lstinline|:=| is the mathematical statement, while the portion after \lstinline|:= by| enters ``tactic mode'' to construct the proof.

The keyword \lstinline|axiom| introduces a statement assumed to be true without providing a proof:

\begin{lstlisting}[language=Lean]
axiom my_axiom : P
\end{lstlisting}

Axioms instruct Lean's kernel to accept \lstinline|P| unconditionally and must be used sparingly to avoid compromising the soundness of the logical system.

To bundle multiple pieces of data together into a single mathematical object or record, Lean uses the \lstinline|structure| keyword.

\begin{lstlisting}[language=Lean]
structure Point where
  x : Float
  y : Float
\end{lstlisting}

This declaration creates a new type called \lstinline|Point|. It automatically generates a
constructor to build points and projection functions (\lstinline|Point.x| and \lstinline|Point.y|) to access the underlying fields.

Finally, Lean handles algebraic hierarchies and overloading via the
\lstinline|class| and \lstinline|instance| keywords. A \lstinline|class| is
essentially a \lstinline|structure| that registers itself with Lean's typeclass
synthesis mechanism (which resolves the square brackets \lstinline|[]| discussed in
Section~\ref{sec:brackets}).

\begin{lstlisting}[language=Lean]
class Add (α : Type) where
  add : α → α → α
\end{lstlisting}

This defines an interface requiring an addition operation for some type \lstinline|α|. To fulfill this requirement for a specific type, we use the \lstinline|instance| keyword to provide the concrete implementation:

\begin{lstlisting}[language=Lean]
instance : Add Nat where
  add x y := x + y
\end{lstlisting}

Once this instance is declared, any function requiring an \lstinline|Add| instance for \lstinline|Nat| will automatically find and use this implementation.

\subsection{The meaning of different brackets}
\label{sec:brackets}

Lean utilizes different types of brackets to distinguish how arguments should be provided or inferred by the system.

Parentheses \lstinline|()| denote \emph{explicit arguments}. These are arguments that the user must manually provide when calling a function or applying a theorem.

\begin{lstlisting}[language=Lean]
def double (n : Nat) : Nat := n + n
\end{lstlisting}

In this example, \lstinline|(n : Nat)| signifies that \lstinline|n| is an explicit argument. Whenever \lstinline|double| is used, the user must explicitly pass a natural number.

Curly brackets \lstinline|{}| denote \emph{implicit arguments}. Lean attempts to infer these arguments automatically from the surrounding context using unification.

\begin{lstlisting}[language=Lean]
def identity {α : Type} (x : α) : α := x
\end{lstlisting}

Here, \lstinline|{α : Type}| indicates that \lstinline|α| is an implicit type parameter. If we apply \lstinline|identity| to a specific natural number, Lean automatically infers that \lstinline|α| must be \lstinline|Nat|, saving the user from writing redundant types.

Square brackets \lstinline|[]| are designated for \emph{typeclass arguments}. While also inferred automatically, they rely on Lean's typeclass synthesis mechanism rather than simple unification.

\begin{lstlisting}[language=Lean]
def addSelf [Add α] (x : α) : α := x + x
\end{lstlisting}

The expression \lstinline|[Add α]| specifies that Lean must find a registered addition structure (a typeclass instance) for the type \lstinline|α|. Lean searches its database to satisfy this requirement automatically when the function is invoked.

\subsection{Attribute tags}
\label{sec: tag_ext_simp}

Lean supports attribute tags, which append special metadata to declarations
to guide automated tools. An attribute is declared directly above a
definition, theorem, or structure using square brackets preceded by an
\lstinline|@| symbol.

The \lstinline|simp| tactic invokes Lean's simplification engine, which applies
a large library of pre-registered lemmas to reduce a goal, often solving
straightforward goals automatically.

A widely used attribute is \lstinline|@[simp]|. Tagging a theorem with \lstinline|@[simp]| adds it to this simplification database, allowing the \lstinline|simp| tactic to apply it autonomously in future proofs.

\begin{lstlisting}[language=Lean]
@[simp]
theorem add_zero (a : Nat) : a + 0 = a := by
  rw [Nat.add_zero]
\end{lstlisting}

Once marked, whenever the simplifier encounters an expression matching \lstinline|a + 0|, it will automatically rewrite it to \lstinline|a| without requiring explicit user instruction.

Another essential attribute is \lstinline|@[ext]|, which establishes extensionality principles. Extensionality asserts that two objects are equal whenever all their corresponding components are equal. Applying \lstinline|@[ext]| directly to a \lstinline|structure| automatically generates the corresponding extensionality lemma.

\begin{lstlisting}[language=Lean]
@[ext]
structure Point where
  x : Float
  y : Float
\end{lstlisting}

Tagging the \lstinline|Point| structure with \lstinline|@[ext]| instructs Lean's metaprogramming framework to generate a theorem behind the scenes. Written manually, it would look like this:

\begin{lstlisting}[language=Lean]
theorem Point.ext {p q : Point} (h₁ : p.x = q.x) (h₂ : p.y = q.y) : p = q
\end{lstlisting}

The \lstinline|@[ext]| tag registers this logic so that invoking the \lstinline|ext| tactic on a goal of \lstinline|p = q| will automatically apply \lstinline|Point.ext|, splitting the goal into the simpler subgoals \lstinline|p.x = q.x| and \lstinline|p.y = q.y|.

Attributes are indispensable for scaling mathematical libraries.
Rather than manually proving and referencing hundreds of structural lemmas
by name, developers register them with tags like \lstinline|@[simp]| and
\lstinline|@[ext]|, enabling Lean's automation to carry the burden in complex
proofs.

\subsection{Namespaces}
\label{sec: namespace}

% \subsection{Namespaces and the \lstinline|open| Keyword}

As mathematical libraries grow, multiple theorems or definitions with similar names become common. Lean uses namespaces to organize declarations and prevent such naming conflicts. When working extensively within a specific context, such as multivariate power series, repeatedly typing the full prefix \lstinline|MvPowerSeries| becomes tedious and clutters the code. Enclosing our code in a \lstinline|namespace| block lets Lean resolve local names automatically, so we may omit the prefix for any definition or theorem belonging to that namespace.

For instance, consider the following explicit definition of a variable $X$:

\begin{lstlisting}[language=Lean]
def MvPowerSeries.X (s : σ) : MvPowerSeries σ R :=
    (MvPowerSeries.monomial (Finsupp.single s 1)) 1
\end{lstlisting}

If we declare that we are operating inside the \lstinline|MvPowerSeries| namespace, the code becomes significantly cleaner. We can drop the prefix from both the new definition's name (\lstinline|X|) and any internal references to existing functions in that namespace (such as \lstinline|monomial|):

\begin{lstlisting}[language=Lean]
namespace MvPowerSeries

def X (s : σ) : MvPowerSeries σ R :=
    monomial (Finsupp.single s 1) 1

end MvPowerSeries
\end{lstlisting}

Once outside the \lstinline|end MvPowerSeries| boundary, a user would once again need to refer to this definition by its fully qualified name: \lstinline|MvPowerSeries.X|.

A complementary tool for managing scopes is the \lstinline|open| keyword.
While \lstinline|namespace| is typically used when \textit{defining} new
objects within a specific theory, \lstinline|open| is used when we simply
want to \textit{access} existing results from another namespace without
typing its prefix. For example, suppose we have a definition that relies on
\lstinline|Finsupp.single|. If we want to use results from the
finitely supported functions library frequently, we can open the
\lstinline|Finsupp| namespace directly:

\begin{lstlisting}[language=Lean]
open Finsupp
namespace MvPowerSeries

def X (s : σ) : MvPowerSeries σ R :=
    monomial (single s 1) 1

end MvPowerSeries
\end{lstlisting}

By writing \lstinline|open Finsupp|, we instruct Lean to look inside the \lstinline|Finsupp| namespace during name resolution, allowing us to reduce \lstinline|Finsupp.single| to just \lstinline|single|.

% \begin{remark}
% The specific functions presented in this example, such as
% \lstinline|monomial| and \lstinline|single|, are utilized here purely
% to illustrate Lean's syntax and scoping mechanisms. We will formally discuss
% the detailed mathematical API for each of these components in Section~\ref{sec:mvpowerseries}.
% \end{remark}


% As mathematical libraries grow, it is common to have multiple theorems or
% definitions with similar names. Lean uses namespaces to organize declarations
% and prevent these naming conflicts. When working extensively within a
% specific mathematical context, such as multivariate power series,
% repeatedly typing the full prefix \lstinline|MvPowerSeries| can become
% tedious and clutter the code. By encapsulating our code within a
% \lstinline|namespace| block, Lean automatically resolves local
% names, allowing us to omit the prefix for any definitions or theorems
% belonging to that namespace.

% For instance, consider the following explicit definition of a variable $X$:

% \begin{lstlisting}[language=Lean]
% def MvPowerSeries.X (s : σ) : MvPowerSeries σ R :=
%     (MvPowerSeries.monomial (Finsupp.single s 1)) 1
% \end{lstlisting}

% If we declare that we are operating inside the \lstinline|MvPowerSeries| namespace, the code becomes significantly cleaner. We can drop the prefix from both the new definition's name (\lstinline|X|) and any internal references to existing functions in that namespace (such as \lstinline|monomial|):

% \begin{lstlisting}[language=Lean]
% namespace MvPowerSeries

% def X (s : σ) : MvPowerSeries σ R :=
%     monomial (Finsupp.single s 1) 1

% end MvPowerSeries
% \end{lstlisting}

% Once outside the \lstinline|end MvPowerSeries| boundary, a user would once again need to refer to this definition by its full, fully qualified name: \lstinline|MvPowerSeries.X|. This mechanism keeps the global namespace clean while ensuring that local development remains readable and concise.

%%%%%%%%%%%%%%%%%%%%%%%%%%%%%%%

% \section{Inductive type and Quotient type}

% \subsection{Inductive type}

% An \emph{inductive type} is defined by specifying its universe level and a comprehensive list of its
% \textit{constructors}. Mathematically, inductive types are ``freely generated''
% by these constructors. Here are several examples in Mathlib

% \begin{lstlisting}
% inductive Nat where
%   | zero : Nat
%   | succ (n : Nat) : Nat

% inductive List (α : Type u) where
%   | nil : List α
%   | cons (head : α) (tail : List α) : List α

% inductive Option (α : Type u) where
%   /-- No value. -/
%   | none : Option α
%   /-- Some value of type `α`. -/
%   | some (val : α) : Option α
% \end{lstlisting}

% In the inductive construction of \lstinline|List|, \lstinline|nil| means the empty list,
% usually written \lstinline|[]|. And the inductive construction of \lstinline|List| is that:
% the list whose first element is \lstinline|head|, where \lstinline|tail| is the rest of the list.
% Usually written \lstinline|head :: tail|.

% Because inductive types are freely generated, Lean automatically enforces two strict rules, often
% referred to as the ``no confusion'' properties \cite{deMoura2021}:
% \begin{enumerate}
%     \item Injectivity of constructors: If $\texttt{succ}(a) = \texttt{succ}(b)$, then $a = b$.
%     \item Disjointness of constructors: Different constructors can never produce the same value.
%     For example, $\texttt{zero} \neq \texttt{succ}(a)$ for any $a$.
% \end{enumerate}

% These properties are what allow functions to be defined via pattern matching and recursive
% evaluation. Is each type in Lean an inductive type? The answer is almost but not.
% Every type in Lean is either an inductive type, or defined in terms of type
% universes (like the type \lstinline|Type 2|), dependent arrow types
% and quotient types.

% \subsection{Quotient types}
% In abstract algebra and topology, it is frequently necessary to construct a new
% space by ``gluing'' elements of a base space together according to an equivalence relation $\sim$.

% If one were to attempt to define a quotient type as a standard inductive type, they would
% face a contradiction. One might try to provide a constructor
% $\texttt{mk} : \alpha \to \text{Quotient}$, but because inductive constructors are strictly
% injective, Lean would enforce that if $a \neq b$, then $\texttt{mk}(a) \neq \texttt{mk}(b)$.
% This entirely defeats the purpose of a quotient, where we explicitly want
% $\texttt{mk}(a) = \texttt{mk}(b)$ whenever $a \sim b$.

% To resolve this, Lean does not treat quotients as inductive types.
% Instead, it extends its core type theory with a set of primitive axioms
% specifically designed for quotienting \cite{Carneiro2019}.

% Lean introduces four foundational constants to handle quotients for any type
% $\alpha$ and relation $r : \alpha \to \alpha \to \texttt{Prop}$:

% \begin{enumerate}
%     \item \lstinline|Quot| is the built-in quotient type.
%     \item \lstinline|Quot.mk| places elements of the underlying type \lstinline|α| into the quotient.
%     \item \lstinline|Quot.lift| allows the definition of functions from the quotient to some other type.
%     \item \lstinline|Quot.ind| is used to write proofs about quotients by assuming that all elements are constructed
%         with \lstinline|Quot.mk|.
% \end{enumerate}

% \begin{lstlisting}
% axiom Quot.sound : ∀ {α : Sort u} {r : α → α → Prop} {a b : α},
%     r a b → Quot.mk r a = Quot.mk r b
% \end{lstlisting}

% By utilizing \texttt{Quot.sound}, Lean intentionally breaks the injectivity rule that governs inductive types, allowing for true mathematical equivalence classes. Functions out of quotient types are then defined using \texttt{Quot.lift}, which requires the user to provide a proof that the function respects the underlying equivalence relation (i.e., the function is well-defined).




%%%%%%%%%%%%%%%%%%%%%%%%%%%%%%%%%%%



% \section{Mathlib: The Mathematical Library}
% \label{sec:mathlib}

% The power of an interactive theorem prover depends heavily on the extent of its standard library. For Lean, this is \textsf{mathlib}, a massive, community-driven repository of formalized mathematics.

% Unlike other theorem provers that often split theories into disconnected libraries, \textsf{mathlib} is heavily integrated and monolithic. This design choice ensures that different areas of mathematics can interact seamlessly. A theorem proven in topology can be immediately applied to algebraic geometry without needing to bridge incompatible definitions of fundamental structures.

% \textsf{mathlib} covers a vast array of topics, including but not limited to:
% \begin{itemize}
%     \item Abstract algebra (groups, rings, modules, category theory)
%     \item Analysis (topology, measure theory, functional analysis)
%     \item Geometry (manifolds, schemes)
%     \item Combinatorics and number theory
% \end{itemize}

% Furthermore, \textsf{mathlib} relies heavily on Lean's typeclass inference system. This system allows Lean to automatically synthesize required mathematical structures—for instance, automatically deducing that the product of two topological spaces is a topological space, or identifying the correct multiplicative inverse operation for an element in a ring.

% By building upon the rigorous foundations of Lean 4 and the extensive mathematical scaffolding of \textsf{mathlib}, this thesis contributes to the ongoing effort to digitize modern mathematics.

% \section{Structures, Classes, and Type Class Inference}
% \label{sec:structures_classes_inference}

% To effectively formalize mathematics, a system must be able to organize abstract concepts—such
% as groups, rings, and topological spaces—in a modular and extensible way. Lean achieves this
% using structures and type classes.

% \subsection{Structures and Classes}
% In Lean, a \lstinline|structure| is an inductive type with a single constructor, akin to a record
% in traditional programming languages. It allows users to bundle data (such as mathematical
% operations) and propositions (such as axioms) into a single cohesive object. For example,
% a group can be defined as a structure containing a carrier type, a multiplication operation,
% an identity element, an inverse operation, and the explicit proofs that these operations
% satisfy the standard group axioms.

% A \lstinline|class| is simply a \lstinline|structure| that has been registered with Lean's type
% class inference system. In mathematical formalization, we almost exclusively use classes
% to define algebraic and topological properties. By doing so, we define generic interfaces
% that many different specific types (e.g., the integers, polynomials, or matrices) can instantiate.

% \subsection{Type Class Inference}
% Type class inference is the automated mechanism by which Lean synthesizes missing
% mathematical context. When a theorem or definition requires an instance of a class—denoted by
% square brackets in Lean's syntax, such as \lstinline|[Ring R]|—Lean searches its
% environment to automatically construct and inject the required instance.

% This system is the absolute backbone of mathematical notation and generic reasoning in Lean.
% For instance, the addition symbol \lstinline|+| is defined via a fundamental type class called
% \lstinline|Add|. When Lean encounters the expression \lstinline|a + b| for variables
% \lstinline|a b : R|, it implicitly triggers type class inference to find an instance of
% \lstinline|Add R|. If \lstinline|R| is known to be a ring via a local assumption,
% Lean seamlessly infers the correct addition operation without the user needing to
% manually specify it. This mechanism allows formalized mathematics to remain readable
% and closely resemble traditional pen-and-paper mathematics.

% \subsection{The Algebraic Hierarchy}
% In \textsf{mathlib}, mathematical structures are organized into a massive, heavily intertwined
% algebraic hierarchy. This hierarchy is built using class extension.
% \begin{lstlisting}
% class Semiring (α : Type u) extends NonUnitalSemiring α, NonAssocSemiring α, MonoidWithZero α

% class Ring (R : Type u) extends Semiring R, AddCommGroup R, AddGroupWithOne R

% class CommRing (α : Type u) extends Ring α, CommMonoid α
% \end{lstlisting}

% For instance,
% the \lstinline|CommRing| (commutative ring) class extends the \lstinline|Ring| class, which in
% turn extends both \lstinline|Semiring|, \lstinline|AddCommGroup| and \lstinline|AddGroupWithOne|,
% cascading all the way down to basic operational classes like \lstinline|Mul| and \lstinline|Add|.

% Because of this hierarchy, a theorem proven for a \lstinline|Monoid| is automatically applicable
% to a \lstinline|Group|, a \lstinline|Ring|, or a \lstinline|Field|. Type class inference handles
% this ``forgetful inheritance'' automatically. When a theorem requires a \lstinline|Monoid| but is
% given a \lstinline|Ring|, Lean's type class synthesis systematically the hierarchy to
% extract the underlying \lstinline|Monoid| structure.

% Designing and maintaining this hierarchy is one of the most significant software engineering
% achievements within the \textsf{mathlib} project. It ensures maximal theorem reuse, preventing
% mathematicians from having to prove the same fundamental lemmas multiple times across different
% domains of algebra.

\section{Structures, classes, and type class inference}
\label{sec:structures_classes_inference}

To formalize mathematics effectively, Lean needs a way to organize
abstract mathematical objects such as groups, rings, modules, and
topological spaces.  These objects usually consist of both data, such as
operations, and properties, such as associativity or distributivity.
Lean organizes this information using \lstinline|structure|,
\lstinline|class|, and \lstinline|instance|.

\subsection{Structures and classes}

A \lstinline|structure| in Lean is similar to a record in programming.
It bundles several pieces of data into a single object.  In mathematics,
this is useful because many objects are naturally described by a
collection of operations together with axioms.  For example, a group
structure consists of a multiplication operation, an identity element, an
inverse operation, and proofs of the group axioms.

A \lstinline|class| is a special kind of structure which is registered
with Lean's type class inference system.  The fields of a class are still
ordinary fields, just as in a structure, but Lean is allowed to search for
instances of the class automatically.  For this reason, algebraic and
topological structures in \textsf{mathlib}, such as \lstinline|Group|,
\lstinline|Ring|, \lstinline|CommRing|, and \lstinline|TopologicalSpace|,
are usually implemented as classes.

For example, a theorem may be stated for an arbitrary type
\lstinline|R| together with an assumption \lstinline|[CommRing R]|:
\begin{lstlisting}
variable (R : Type*) [CommRing R]
\end{lstlisting}
The square brackets mean that Lean should treat \lstinline|CommRing R| as
an instance argument.  In practice, this means that the ring operations
and ring axioms on \lstinline|R| are available automatically.

This is different from passing an ordinary proposition explicitly.  A
class is appropriate when the structure should be found automatically and
used throughout many definitions and theorems.  By contrast, an ordinary
proposition is more suitable when the assumption is local and should be
passed explicitly.

\subsection{Instances}

An \lstinline|instance| provides a particular realization of a class for a
specific type.  For example, \textsf{mathlib} contains instances showing
that the integers form a ring, that polynomial rings inherit ring
structures, and that many constructions preserve algebraic structures.

Once an instance is available, Lean can use it without the user having to
name it explicitly.  For instance, if Lean knows \lstinline|[CommRing R]|,
then it can use the addition, multiplication, zero, one, negation, and all
the ring laws associated with \lstinline|R|.  This is why ordinary
mathematical notation such as \lstinline|x + y|, \lstinline|x * y|, and
\lstinline|-x| can be used in a general ring.

This mechanism is also important in the formalization of formal group
laws.  For example, after constructing an \lstinline|AddCommGroup|
instance on \lstinline|F.Point σ|, Lean can use the notation
\lstinline|x + y|, \lstinline|-x|, and \lstinline|n • x| for points of a
formal group law without any additional manual specification.

\subsection{Type class inference}

Type class inference is the mechanism by which Lean automatically finds
the required instances.  When a definition or theorem requires an
assumption such as \lstinline|[Ring R]| or \lstinline|[TopologicalSpace X]|,
Lean searches its environment for an appropriate instance and inserts it
automatically.

For example, the notation \lstinline|a + b| requires Lean to know an
addition operation on the type of \lstinline|a| and \lstinline|b|.  If
\lstinline|a b : R| and Lean has an instance \lstinline|[Ring R]|, then
Lean can infer the required addition operation from the ring structure.
Thus the user can write expressions in a style close to ordinary
mathematics, while Lean fills in the underlying algebraic data.

Type class inference is especially useful for generic mathematics.  A
single theorem can be stated for all rings, groups, or modules, and then
applied to any concrete type for which Lean can synthesize the required
instances.

\subsection{The Algebraic Hierarchy}

In \textsf{mathlib}, algebraic structures are organized into a large
hierarchy.  More complicated structures extend simpler ones.  For example,
a ring contains an additive commutative group structure and a
multiplicative monoid structure, together with distributivity laws.
A commutative ring further extends a ring by adding commutativity of
multiplication.

This hierarchy is implemented using class extension.  For example:
\begin{lstlisting}
class Semiring (α : Type u) extends NonUnitalSemiring α,
  NonAssocSemiring α, MonoidWithZero α

class Ring (R : Type u) extends Semiring R, AddCommGroup R,
  AddGroupWithOne R

class CommRing (α : Type u) extends Ring α, CommMonoid α
\end{lstlisting}

Because of this hierarchy, an instance of \lstinline|CommRing R| also
provides access to the structures inherited from its parent classes, such
as \lstinline|Ring R|, \lstinline|AddCommGroup R|, and the underlying
addition and multiplication operations.  Therefore, a theorem proved for a
more general structure can automatically be applied to a more specific
one.  For instance, a theorem about additive commutative groups can be
used for any commutative ring, because a commutative ring contains an
additive commutative group structure.

This design is essential for theorem reuse in \textsf{mathlib}.  Instead
of reproving the same basic facts separately for groups, rings, fields,
and other structures, one proves a theorem at the appropriate level of
generality, and type class inference makes it available in all more
specialized contexts.